\section{Condiciones de las realizaciones}
\label{sec:ii-condiciones-realizaciones}

Las realizaciones de acción, mezcla, área y gravitación conservan
los dominios y las normalizaciones de los resultados que las componen.

\subsection{Coordenada de alfa y marcos de mezcla}

La clausura dodecafásica y la representación analítica de alfa expresan
la misma coordenada a toda profundidad
(teorema~\ref{thm:alpha-dos-publicaciones-completas}).
La recurrencia analítica se fija desde el estado enriquecido y la clase
geométrica declarada; no selecciona independientemente esa coordenada.
En la realización angular, los marcos espectrales completos intervienen
en \eqref{eq:ii-pmns-frames}. La unitariedad de un factor polar no sustituye
su identificación con ese cambio de marcos. La compatibilidad CP conserva
la condición \eqref{eq:ii-cp-square}, y una forma real de Majorana restringe
las fases admisibles.

\subsection{Bases térmicas y normalización gravitatoria}

La sección~\ref{sec:ii-boltzmann} define el transductor térmico y su
transformación bajo cambio de bases. El producto
\(k_B\Theta_{\rm clk}\log3\) queda fijado por la energía del ciclo;
la coordenada individual de \(k_B\) exige especificar las bases positivas
de las rectas de energía y temperatura. Esta covariancia no altera
la acción ni el período precedentes.

La sección~\ref{sec:gravity} fija la coordenatización gravitatoria mediante el radio
reducido \(Gm/c^2\) y su condición de coincidencia con \(R_{108}\).
En la composición actual, la longitud precede a esa expresión circular:
R4 y R5 determinan \(L\) y el radio positivo del carácter, y el reloj
\(t_0=\pi_{\HMT}L/(54c_{\rm int})\) da \(R_{108}=L\).
R8 identifica este radio con la lectura gravitatoria reducida.
Las secciones de acción recíprocas conservan el producto \(\hbar G\)
y, por ello, la escala areal correspondiente. La ecuación de Cartan--Holst
se formula para una cotétrada no degenerada, una conexión y la corriente
de espín declaradas, con las condiciones de contorno de su variación.

\subsection{Procedencia de la composición longitudinal y radial}

La conservación del archivo, el transporte de Hadamard, la acción de
fase y la reciprocidad de las secciones proceden de los desarrollos
anteriores del corpus. La revisión del 25--26 de septiembre de 2026
reúne el registro marcado de incidencias y explicita las composiciones
longitudinal R4 y métrica R5. El teorema de
\ref{g26:sec:rigidez-radial-es} demuestra su consecuencia conjunta con
R1--R3 y R6--R8: unicidad de la inscripción, del lector radial y del
coeficiente común en la realización declarada \cite{hmtg26radial}.
Esta procedencia distingue los antecedentes conservados de las
definiciones y pruebas incorporadas en la revisión.

Los transportes con métrica global, incluidos sus bloques cruzados,
conservan la acción y el residuo recuperable. La restitución del
contraángulo y de la velocidad constitutiva proporciona expresiones
equivalentes de la misma evaluación \cite{hmtg26metric,hmtg26counter}.
La sección longitudinal de referencia, la acción y la velocidad
dimensional permanecen explícitas. La elección \(L_*=1\,\mathrm m\)
es la sección utilizada para las cifras SI; su declaración permite
reproducirlas con el mismo significado dimensional. La identificación
física R8 conserva su lugar en la demostración.
